\documentclass[11pt,a4paper]{article}
\usepackage[utf8]{inputenc}
\usepackage[margin=1in]{geometry}
\usepackage{graphicx}
\usepackage{hyperref}
\usepackage{listings}
\usepackage{xcolor}
\usepackage{enumitem}
\usepackage{fancyhdr}
\usepackage{tabularx}
\usepackage{array}
\usepackage{booktabs}
\usepackage{amssymb}
\usepackage{tikz}

% Define colors for code listings
\definecolor{codegreen}{rgb}{0,0.6,0}
\definecolor{codegray}{rgb}{0.5,0.5,0.5}
\definecolor{codepurple}{rgb}{0.58,0,0.82}
\definecolor{backcolour}{rgb}{0.95,0.95,0.92}
\definecolor{goldcolor}{rgb}{0.85,0.65,0.13}

% Code listing style
\lstdefinestyle{mystyle}{
    backgroundcolor=\color{backcolour},
    commentstyle=\color{codegreen},
    keywordstyle=\color{magenta},
    numberstyle=\tiny\color{codegray},
    stringstyle=\color{codepurple},
    basicstyle=\ttfamily\footnotesize,
    breakatwhitespace=false,
    breaklines=true,
    captionpos=b,
    keepspaces=true,
    numbers=left,
    numbersep=5pt,
    showspaces=false,
    showstringspaces=false,
    showtabs=false,
    tabsize=2,
    frame=single
}

\lstset{style=mystyle}

% Page style
\pagestyle{fancy}
\fancyhf{}
\rhead{2025-08-26}
\lhead{submission\_instructions\_v1.md}
\cfoot{\thepage\ / 7}
\renewcommand{\headrulewidth}{0.4pt}
\renewcommand{\footrulewidth}{0.4pt}

\title{\textbf{TechArena 2025 - Phase 1 Submission Instructions}}
\author{}
\date{Last Updated: 26-08-2025}

\begin{document}

\maketitle
\thispagestyle{fancy}

\section*{📋 Overview}

Welcome to TechArena 2025! This document provides comprehensive instructions for submitting your Phase 1 solution. Please read all sections carefully to ensure your submission is properly validated and evaluated.

\section*{🎯 Submission Requirements}

\subsection*{Required Deliverables}

Your submission must include \textbf{all} of the following components:

\begin{enumerate}
    \item \textbf{Main Script}: \texttt{main.py} (in project root directory)
    \item \textbf{Input Data Processing}: Code to read and process the provided Excel input file
    \item \textbf{Output Files}: Three result files in CSV or Excel format
    \item \textbf{Dependencies}: \texttt{requirements.txt} file (if using external packages)
    \item \textbf{Documentation}: \texttt{README.md} file explaining your approach
\end{enumerate}

\subsection*{Output Files Required}

Your \texttt{main.py} script must generate exactly \textbf{three output files} with the following names (CSV or Excel format):

\begin{enumerate}
    \item \textbf{TechArena\_Phase1\_Configuration.csv} (or \texttt{.xlsx})
    \begin{itemize}[leftmargin=*]
        \item Optimal BESS configuration parameters
        \item Required columns: C-rate, number of cycles, yearly profits [kEUR/MW], levelized ROI [\%]
    \end{itemize}
    
    \item \textbf{TechArena\_Phase1\_Investment.csv} (or \texttt{.xlsx})
    \begin{itemize}[leftmargin=*]
        \item Investment analysis and ROI calculations
        \item Must include: WACC, inflation rate, discount rate, yearly profits, year-by-year analysis, levelized ROI
    \end{itemize}
    
    \item \textbf{TechArena\_Phase1\_Operation.csv} (or \texttt{.xlsx})
    \begin{itemize}[leftmargin=*]
        \item Required columns: Timestamp, Stored energy [MWh], SoC [-], Charge [MWh], Discharge [MWh], Day-ahead buy [MWh], Day-ahead sell [MWh], FCR Capacity [MW], aFRR Capacity POS [MW], aFRR Capacity NEG [MW]
    \end{itemize}
\end{enumerate}

\section*{📁 Project Structure}

Your submission should follow this \textbf{exact} directory structure:

\begin{lstlisting}[language=bash, caption=Required Directory Structure]
your_team_submission.zip
├── main.py                 # Main execution script (REQUIRED)
├── requirements.txt        # Python dependencies (if needed)
├── README.md              # Documentation (RECOMMENDED)
├── input/                 # Input data directory
│   └── TechArena2025_ElectricityPriceData_v2.xlsx
└── [additional files/folders]  # Your implementation files
    ├── output/            # Output data directory
\end{lstlisting}

\section*{🚨 Critical Requirements}

\begin{itemize}
    \item \texttt{main.py} \textbf{MUST be in the root directory} (not in subdirectories)
    \item Input file \textbf{MUST be in} \texttt{input/} subdirectory
    \item Output files will be generated in the output directory
    \item \textbf{No absolute file paths} - use relative paths only
\end{itemize}

\section*{💻 Technical Specifications}

\subsection*{Python Requirements}

\begin{itemize}
    \item \textbf{Python Version}: 3.8 or higher
    \item \textbf{Output Format}: CSV (.csv) or Excel (.xlsx) files
\end{itemize}

\subsection*{Input Data Structure}

The input file contains 4 sheets:

\begin{itemize}
    \item \textbf{Day-ahead prices}: Historical electricity prices for 2024
    \item \textbf{FCR prices}: Frequency Containment Reserve prices
    \item \textbf{aFRR capacity prices}: Automatic Frequency Restoration Reserve prices
    \item \textbf{Data description}: Market conditions and BESS parameters
\end{itemize}

\subsection*{Data Validation Constraints}

\textbf{Configuration File:}
\begin{itemize}
    \item Must include required parameters: yearly profits [kEUR/MW], levelized ROI [\%] for all C-rate and number of cycles combinations
\end{itemize}

\textbf{Investment File:}
\begin{itemize}
    \item Must include required parameters: WACC, inflation rate, discount rate
    \item Must include Levelized ROI calculation
    \item Below is the sample table for reference:
\end{itemize}

\begin{table}[h!]
\centering
\begin{tabular}{|l|c|}
\hline
\textbf{Parameter} & \textbf{Value} \\
\hline
WACC & value \\
\hline
Inflation Rate & value \\
\hline
Discount rate & value \\
\hline
Yearly profits (2024) & value \\
\hline
\end{tabular}
\end{table}

\begin{table}[h!]
\centering
\begin{tabular}{|c|c|c|}
\hline
\textbf{Year} & \textbf{Initial Investment [kEUR/MWh]} & \textbf{Yearly profits [kEUR/MWh]} \\
\hline
2023 & & \\
2024 & & \\
2025 & & \\
2026 & & \\
2027 & & \\
2028 & & \\
2029 & & \\
2030 & & \\
2031 & & \\
2032 & & \\
2033 & & \\
\hline
\multicolumn{2}{|l|}{Levelized ROI} & Value \\
\hline
\end{tabular}
\end{table}

\textbf{Operation File:}
\begin{itemize}
    \item Must include required parameters: Timestamp, Stored energy [MWh], SoC [-], Charge [MWh], Discharge [MWh], Day-ahead buy [MWh], Day-ahead sell [MWh], FCR Capacity [MW], aFRR Capacity POS [MW], aFRR Capacity NEG [MW]
\end{itemize}

\section*{🔧 Implementation Guidelines}

\subsection*{Example main.py Structure}

\begin{lstlisting}[language=Python, caption=Example main.py Structure]
import os
import pandas as pd

def main():
    """
    Main execution function for TechArena Phase 1 solution
    """
    # Read input data
    input_file = os.path.join("input", "TechArena2025_ElectricityPriceData.xlsx")
    
    # Process data and run optimization
    # ... your implementation here ...
    
    # Generate output files
    config_df.to_csv("TechArena_Phase1_Configuration.csv", index=False)
    investment_df.to_csv("TechArena_Phase1_Investment.csv", index=False)
    operation_df.to_csv("TechArena_Phase1_Operation.csv", index=False)
    
    print("✅ All output files generated successfully!")

if __name__ == "__main__":
    main()
\end{lstlisting}

\subsection*{Best Practices}

\textbf{✅ DO:}
\begin{itemize}
    \item Use relative paths for all file operations
    \item Check if input files exist before processing
    \item Handle exceptions gracefully with try-catch blocks
    \item Validate your output before submission
    \item Test your script in a clean environment
    \item Use meaningful variable names and add comments
    \item Include all required dependencies in requirements.txt
\end{itemize}

\textbf{❌ DON'T:}
\begin{itemize}
    \item Use absolute file paths (e.g., \texttt{C:\textbackslash Users\textbackslash YourName\textbackslash...})
    \item Hardcode team-specific paths or credentials
    \item Rely on files not provided in the input
    \item Create infinite loops or excessive processing
    \item Leave debugging print statements that flood output
    \item Submit incomplete or placeholder files
\end{itemize}

\subsection*{Sample requirements.txt}

\begin{lstlisting}[caption=Sample requirements.txt]
pandas>=1.5.0
numpy>=1.21.0
openpyxl>=3.0.9
matplotlib>=3.5.0
scipy>=1.8.0
\end{lstlisting}

\section*{📝 Submission Process}

\subsection*{Step 1: Prepare Your Submission}

\begin{enumerate}
    \item Complete your implementation in \texttt{main.py}
    \item Test thoroughly with the provided input data
    \item Verify all output files are generated correctly
    \item Create/update \texttt{requirements.txt} with all dependencies
    \item Write \texttt{README.md} explaining your approach and any special instructions
\end{enumerate}

\subsection*{Step 2: Test Your Submission}

Before submitting, test your solution:

\begin{lstlisting}[language=bash, caption=Testing Commands]
# Create a clean test environment
cd your_project_directory

# Install dependencies (if using virtual environment)
pip install -r requirements.txt

# Run your main script
python main.py

# Verify output files are created
ls -la *.csv  # or *.xlsx
\end{lstlisting}

\subsection*{Step 3: Create Submission Package}

\begin{enumerate}
    \item Create a ZIP file containing your entire project
    \item Name the file: \texttt{TeamName\_TechArena2025\_Phase1.zip}
    \item Verify contents match the required structure
    \item Test the ZIP file by extracting and running in a new location
\end{enumerate}

\subsection*{Step 4: Submit}

\begin{enumerate}
    \item Upload your ZIP file to the designated submission platform
    \item Confirm successful submission before the deadline
\end{enumerate}

\section*{⚡ Automated Validation}

Your submission will undergo \textbf{automated validation} that checks:

\textbf{✅ Structure Validation}
\begin{itemize}
    \item Correct file and directory structure
    \item Presence of required files (\texttt{main.py}, output files)
    \item Proper file naming conventions
\end{itemize}

\textbf{✅ Code Quality}
\begin{itemize}
    \item Script executability and syntax
    \item Import statement validation
    \item Function and documentation presence
    \item Excel/CSV handling capabilities
\end{itemize}

\textbf{✅ Execution Testing}
\begin{itemize}
    \item Virtual environment setup and dependency installation
    \item Error handling and output generation
\end{itemize}

\textbf{✅ Output Validation}
\begin{itemize}
    \item Correct number and format of output files
    \item Schema compliance (column names and types)
    \item File size and content reasonableness
\end{itemize}

\section*{🚨 Common Issues and Solutions}

\begin{description}
    \item[Issue:] ``main.py not found in project root''
    \item[Solution:] Ensure \texttt{main.py} is directly in the ZIP file root, not in a subdirectory.
    
    \item[Issue:] ``Input file not found''
    \item[Solution:] Place input file in \texttt{input/} subdirectory and use relative path:
\end{description}

\begin{lstlisting}[language=Python]
input_file = os.path.join("input", "TechArena2025_ElectricityPriceData_v2.xlsx")
\end{lstlisting}

\begin{description}
    \item[Issue:] ``Requirements installation failed''
    \item[Solution:] \hfill
    \begin{itemize}
        \item Test your \texttt{requirements.txt} in a clean environment
        \item Use specific version numbers for critical packages
        \item Include \texttt{openpyxl} for Excel file handling
    \end{itemize}
    
    \item[Issue:] ``Output file validation failed''
    \item[Solution:] \hfill
    \begin{itemize}
        \item Check column names match exactly (case-sensitive)
        \item Verify data types and ranges
        \item Ensure no missing values in required fields
    \end{itemize}
    
    \item[Issue:] ``Script execution timeout''
    \item[Solution:] \hfill
    \begin{itemize}
        \item Optimize your algorithms for efficiency
        \item Avoid unnecessary data loading or processing
        \item Consider using vectorized operations with pandas/numpy
    \end{itemize}
\end{description}

\section*{📞 Support and Contact}

\subsection*{Technical Support}
\begin{itemize}
    \item \textbf{Email}: \href{mailto:amandine@bemyapp.com}{amandine@bemyapp.com}
\end{itemize}

\subsection*{Only For Code Submission Support}
\begin{itemize}
    \item \textbf{Email}: \href{mailto:mangesh.mankar@huawei.com}{mangesh.mankar@huawei.com}
\end{itemize}

\subsection*{FAQ and Updates}
\begin{itemize}
    \item \textbf{Event Website}: \url{https://techarena-nuremberg.hackathon.com/}
    \item \textbf{Discussion Forum}: Available on the event platform
    \item \textbf{Live Q\&A Sessions}: Scheduled times will be announced
\end{itemize}

\section*{📋 Submission Checklist}

Before submitting, verify:

\begin{itemize}
    \item[$\square$] \texttt{main.py} is in the root directory
    \item[$\square$] Input file is in \texttt{input/} subdirectory
    \item[$\square$] Script runs successfully and generates all 3 output files
    \item[$\square$] \texttt{requirements.txt} includes all dependencies
    \item[$\square$] README.md explains your approach
    \item[$\square$] ZIP file is named correctly: \texttt{TeamName\_TechArena2025\_Phase1.zip}
    \item[$\square$] No absolute file paths or hardcoded paths
    \item[$\square$] All output files have correct column names
    \item[$\square$] Tested in a clean environment
\end{itemize}

\vspace{1cm}

\begin{center}
\Large{\textbf{Good luck, and may the best solution win! 🚀}}
\end{center}

\vspace{1cm}

\hrule
\vspace{0.5cm}
\textit{TechArena 2025 Organizing Committee}\\
\textit{Last Updated: 26-08-2025}

\end{document}